\chapter{Differentiable molecular and sequence models}
\ChapterGoals{%
\item specify exactly which quantity is differentiated, with respect to which parameters;
\item derive reaction sensitivities and check them by three independent routes;
\item differentiate a discretised solver and detect a graph-breaking implementation bug;
\item relax nucleotide sequences to a probability simplex and compute exact expected scores;
\item recognise when a relaxation optimises a reward that no real sequence achieves, and why design needs independent evidence.}
{Chapters~7 and~28; derivatives and the chain rule; basic PyTorch (introduced as needed).}
Chapter 28 gave a reaction system a clock. Now ask a different question:
if the modeled terminal population misses a target, which parameter should move?
And if a sequence predictor assigns a low score, can a gradient suggest a better
sequence? The answer depends on precisely \emph{what} is differentiated.

This chapter constructs two transparent examples rather than hiding the problem in
a pretrained network. We derive reaction sensitivities, compare three derivative
routes, differentiate a discretized solver, and expose a graph-breaking bug.
Then we derive a nucleotide-logit gradient, enumerate its discrete expectation,
and test a constraint that a continuous relaxation does not preserve.
All computed values are teaching-model outputs, not measurements.

\section{Choose the differentiable object}
A physical molecule, an integer reaction trajectory, a real-valued kinetic model
and a learned sequence predictor are different objects. Write the full contract
before drawing the computational graph:
\[
  \theta\in\Theta,\quad \dot x=F(x,\theta,t),\quad
  y=H(x(T),\theta),\quad J=\mathcal L(y,y^\star).
\]
Topology, initial state, units, observation mapping and objective all belong in
this contract. Automatic differentiation traverses the executed operations, not
the validity of the modeling arrow. A perfect derivative of a poor predictor
can be a poor scientific recommendation.
\Plate{01-boundary}{A reversible reaction is mapped into an explicit mathematical
program. The backward arrows terminate inside that program, not in the physical
vessel. Rates and observation units stay visible.}{dnad29:boundary}

\Term{Sensitivity} is a derivative of a specified output with respect to a specified
parameter. If $J$ has population-squared units and $k_f$ has units s$^{-1}$,
$dJ/dk_f$ has population-squared seconds. Derivative magnitude changes when rate
units or parameter coordinates change. To compare dimensions, introduce reference
scales $x=C\tilde x$, $t=t_0\tilde t$ and $\kappa=k t_0$, then state whether
the reported derivative is with respect to $k$ or $\kappa$.

\section{Derive both reaction sensitivities}
Use the Chapter 28 first-order, well-mixed, constant-rate model
\[
 A\xrightleftharpoons[k_r]{k_f}B,\qquad
 \dot A=-k_f A+k_r B,\quad C=A_0+B_0.
\]
Counts in the deterministic model may be interpreted as ensemble means under
the matching first-order stochastic process; they are not necessarily integers.
Let $\lambda=k_f+k_r>0$, $E=e^{-\lambda t}$,
$q=1-E$ and $D=A_0-Ck_r/\lambda$. Then
\[
 A(t)=Ck_r/\lambda+D E.
\]
Differentiate the equilibrium and transient terms separately:
\begin{align}
 s_f(t):=\partial_{k_f}A(t)&=-\frac{Ck_r}{\lambda^2}q-DtE,\\
 s_r(t):=\partial_{k_r}A(t)&=\frac{Ck_f}{\lambda^2}q-DtE.
\end{align}
The signs have a concrete interpretation for the fixture: increasing forward
conversion lowers $A$; increasing reverse conversion replenishes it.
Both sensitivities are zero at $t=0$ because the stipulated initial populations
do not depend on the rates. A rate-dependent initialization would add its own term.

Set $A_0=100$, $B_0=0$, $k_f=.4$ s$^{-1}$, $k_r=.1$ s$^{-1}$,
$T=5$ s and $A^\star=30$. We obtain
\[
 A(T)=26.56679989,\quad s_f=-69.55059950,\quad s_r=114.03240077.
\]
For $J=(A(T)-A^\star)^2$,
\[
 J=11.786862996,\qquad
 \frac{dJ}{dk_f}=2(A(T)-A^\star)s_f=477.562251752.
\]
A small decrease of $k_f$ is therefore a local descent direction. This is neither
a globally optimal rate nor a calibrated chemical design.
\Plate{02-sensitivity}{Computed concentration/count-like output and rate
sensitivities have separate vertical scales and units. A perturbation is propagated
through both equilibrium and relaxation speed.}{dnad29:sensitivity}
\Listing{analytic_sensitivities}

\section{Implement without breaking the graph}
Float64 is explicit rather than a module-wide change of PyTorch's default dtype.
The following function keeps rate tensors on the graph and evaluates
$1-e^{-\lambda t}$ as $-\operatorname{expm1}(-\lambda t)$.
\Listing{reversible_a_t}
Validation may inspect detached values, but the returned calculation must use the
original tensors. Converting the rate itself to a Python float before evaluation
would erase the claimed dependency. This chart deliberately requires $\lambda>0$.
The all-zero-rate boundary has a separate constant forward solution; extending
derivatives there requires a continuous-limit treatment, not division by zero.

\section{Three routes, one declared target}
Use the analytical formula above, Float64 reverse-mode AD, and central differences
of an independent Python/math implementation:
\[
 g_{\rm FD}(\epsilon)=
 \frac{J_{\rm scalar}(k_f+\epsilon)-J_{\rm scalar}(k_f-\epsilon)}{2\epsilon}.
\]
The finite-difference function does not call the AD loss wrapper or its backward
pass. This independence matters: two wrappers around one wrong derivative do not
provide two oracles.
\begin{center}\input{\Generated/gradient-checks.tex}\end{center}
Large steps incur truncation; extremely small steps amplify subtraction error.
Here $10^{-6}$ is more accurate than $10^{-8}$. Sweep steps and parameters rather
than selecting one favorable result. Reject unrepresentable steps and perturbations
outside the rate domain.
\Plate{03-oracles}{The analytical, independent scalar finite-difference and AD
paths converge on one derivative target. A separate comparison is required for
the derivative of a numerical solver.}{dnad29:oracles}

For $f:\mathbb R^p\to\mathbb R^q$, a Jacobian--vector product propagates one input
direction; a vector--Jacobian product propagates one output cotangent. Reverse mode
is attractive for a scalar loss with many parameters. It computes the derivative
of the selected program trace where differentiable, not a license to ignore branches,
discontinuities or numerical error.

\section{Differentiate the solver you actually executed}
For step size $h=T/n$, forward Euler gives
\[
 A_{j+1}=(1-\lambda h)A_j+h k_r C.
\]
Its derivative with respect to $k_f$, with $k_r,C,h$ fixed, is
\[
 s_{j+1}=(1-\lambda h)s_j-h A_j,\qquad s_0=0.
\]
The $A_j$ on the right is the \emph{old} state. Updating $A$ first and then using
it here implements a different recurrence.
\Listing{euler_sensitivity}
The reference enforces Chapter 28's positivity condition
$h\max(k_f,k_r)\leq1$. A successful gradient calculation does not repair
a negative-population numerical method.

\begin{center}\input{\Generated/euler-checks.tex}\end{center}
AD agrees with the manual recurrence and an independent geometric-series oracle
at each fixed $h$. It does \emph{not} equal the continuous-model derivative at
finite $h$; the gradient error approximately halves under these refinements.
Separate \Term{solver derivative} correctness from convergence to the underlying
differential-equation derivative.

\subsection{An intentionally wrong backward path}
At $h=.1$, repeatedly replacing the previous state by a detached tensor preserves
$A(T)=26.15559802$ exactly. Yet the resulting loss gradient is $20.35962753$,
rather than $532.98854619$. A forward-only parity test cannot catch this mutation.
Compare every claimed parameter gradient, not only the final state.

For larger systems, forward sensitivity obeys
$\dot S=F_xS+F_\theta$. An adjoint instead propagates output sensitivities backward.
Discrete backpropagation differentiates an executed integration scheme;
a continuous adjoint differentiates a continuous model and is itself solved
numerically. Check interpolation, tolerances, stiffness, event handling and
memory/recomputation policies before declaring these interchangeable.

\section{Constraints change coordinates}
A positive rate can be parameterized as $k_f=k_0\operatorname{softplus}(\rho)$,
where $\rho$ is dimensionless and $k_0$ supplies inverse-time units. Then
\[
 \partial_\rho J=(\partial_{k_f}J)\,k_0\,\sigma(\rho).
\]
The reference checks this chain rule at $\rho=-.7$ with $k_0=1$ s$^{-1}$.
Softplus enforces strict positivity but may have a tiny derivative far below zero;
it neither represents exact zero at a finite $\rho$ nor fixes poor conditioning.
A positive parameterization is not an identifiability argument.

\section{Can the data determine both rates?}
\Term{Local identifiability} requires informative observations.
A single terminal observation has sensitivity row $(s_f,s_r)$. The direction
$(s_r,-s_f)$ has zero first-order effect. Thus one scalar observation cannot
locally identify two unconstrained rates. With observations at two times:
\begin{center}\input{\Generated/identifiability.tex}\end{center}
the determinant is $-6324.063041$, nonzero at this fixture.
This gives local rank two, not global uniqueness or small uncertainty.
At equilibrium only $k_r/(k_f+k_r)$ is observed: doubling both rates leaves
the equilibrium unchanged.

With measurement errors $\sigma_i$, use a whitened sensitivity matrix
$\tilde S_{ij}=(\partial_{\theta_j}y_i)/\sigma_i$ and inspect singular directions.
Under the stipulated independent Gaussian observation model, the local
information matrix is $\tilde S^\top\tilde S$. Correlated errors instead require
a covariance-weighted construction. A poorly conditioned matrix warns that
some rate combinations remain weakly constrained. Observation times, initial
conditions and perturbations should be designed to distinguish those combinations.

\section{From nucleotide symbols to a simplex}
A discrete base has no ordinary derivative with respect to its letter.
Introduce logits $\ell\in\mathbb R^{L\times4}$ in A,C,G,T order and temperature
$\tau>0$:
\[
 p_{ib}=\frac{e^{\ell_{ib}/\tau}}{\sum_c e^{\ell_{ic}/\tau}}.
\]
These rows are categorical distributions, not four fractional nucleotides
inside a single DNA molecule. A physical library sampled from them is yet another
object with sampling and preparation costs.
\Plate{04-relaxation}{Four nucleotide channels become a probability simplex,
then a score and a discrete candidate. The discretization arrow has its own
constraint and score checks.}{dnad29:relaxation}

For a fixed motif $m$ of length $L$, define
$S(p)=L^{-1}\sum_i p_{i,m_i}$. It is exactly the expected positional match fraction
under independent base sampling. Enumerating all $4^4=256$ strings provides an
independent expectation oracle. Its logit derivative is
\[
 \frac{\partial S}{\partial\ell_{ib}}
 =\frac{p_{i,m_i}}{L\tau}\left(\mathbf1[b=m_i]-p_{ib}\right).
\]
\Listing{soft_motif_gradient}
At uniform logits, $L=4,\tau=1$, the target-channel gradient is $.046875$;
each other channel has gradient $-.015625$. Each row sums to zero, expressing
the softmax gauge: adding the same constant to all four logits changes nothing.
All sixteen entries are checked against AD across three temperatures.

Eighty plain gradient-ascent updates with step $.4$ raise the score from $.25$
to $.780406667$; the argmax candidate is ACGT and every row remains non-one-hot.
History includes the initial score \emph{and} the score after each update.
This is a known-answer objective, not evidence of promoter, enhancer or protein
design ability.

\section{A relaxation can invent a nonexistent reward}
Consider the nonlinear score
\[
 R(p)=\frac1L\sum_i\left(1-\sum_b p_{ib}^2\right).
\]
Uniform rows achieve $.75$, but \emph{every} discrete one-hot sequence scores zero.
This is a \Term{relaxation gap}: $R(\mathbb E[X])\neq\mathbb E[R(X)]$. Improving a predictor evaluated
on a soft sequence does not necessarily improve the expected discrete sequence
or the eventual argmax candidate.

Likewise, uniform length-four rows have expected GC count two, but A-first
argmax tie-breaking returns AAAA with GC count zero. A differentiable expected-GC
penalty is not an exact per-candidate GC constraint.
\par\medskip\noindent\begin{minipage}{\linewidth}
\Listing{exact_constrained_design}
\end{minipage}\par\medskip
An exhaustive oracle for the AAAA positional target with exactly two GC bases
returns AACC, score $.5$, under lexical tie-breaking. The enumeration is limited
to length eight, costs exponentially many candidates, and is a correctness oracle,
not a production search method.

\section{Discrete and stochastic derivatives need precise names}
For categorical distributions, one may use a score-function estimator
\[
 \nabla_\theta\mathbb E_{X\sim p_\theta}[f(X)]
 =\mathbb E[f(X)\nabla_\theta\log p_\theta(X)]
\]
under the relevant support and interchange conditions. Baselines can reduce
variance; biased straight-through rules replace a discrete forward choice's
backward map by a chosen surrogate. Neither is the ordinary derivative of
argmax. Linder and Seelig's Fast SeqProp combines normalization with a
straight-through approach; its results motivate the distinction between relaxed
and discrete predictor inputs \cite{LinderSeelig2021}. Our reference is not a
Fast SeqProp reproduction.

An exact SSA path changes reaction choices and horizon event counts discontinuously
as rates vary. A locally unchanged integer endpoint can have zero pathwise
derivative while its expected count changes smoothly. Chapter 28's independent-
molecule expectation gives an exact derivative oracle for this first-order case;
it does not supply a differentiable exact SSA implementation.
Rijal and Mehta smooth discrete operations in an approximate differentiable
Gillespie method \cite{RijalMehta2025}. Compare distributions, moments and task-relevant
tails to exact SSA, not only attractive gradients.

\begin{CSBridge}
\emph{Automatic differentiation} records every arithmetic operation of a program
in a graph and applies the chain rule backwards through it, giving exact
derivatives of the computed output with respect to every input. Frameworks such
as PyTorch do this automatically for tensors. The derivative is of the program
that actually ran, including its discretisation and any accidental
\texttt{detach}, which is why this chapter tests gradients against independent
calculations.
\end{CSBridge}

\section{A design loop needs independent evidence}
\Plate{05-design-loop}{Optimization creates proposals. Discretization, constraints,
unused validation evidence and physical testing are separate gates; a failure
returns to the model or design, not to a relabeled success.}{dnad29:design-loop}
Freeze predictor training support and design objectives; optimize only inside a
declared design domain; discretize; recheck exact sequence constraints; evaluate
with evidence not used for candidate selection; then measure physical function
where the claim requires it. A second predictor trained on the same biases may not
be genuinely orthogonal evidence. Record all candidates, exclusions and selection
rules, not only the best predicted score.

Abstract mass-action CRNs have been optimized using differentiable programming
\cite{Mordvintsev2023}. Indexed CRNs provide a contemporary framework for families of
species and reactions, including well-mixed and reaction-diffusion settings
\cite{Lee2025}. These are research connections, not evidence that our two-state
example implements those systems or produces realizable DNA kinetics.

\Plate{06-failures}{Failure signatures distinguish graph errors, solver error,
non-identifiability and relaxation mismatch. Each has a different remedy and
a different falsifying test.}{dnad29:failures}

\begin{KeyIdeas}
\begin{itemize}\setlength{\itemsep}{2pt}
\item A gradient is only meaningful once the model, the parameters, the solver and
the objective are declared.
\item Analytic sensitivities, finite differences and automatic differentiation
should agree on the same target; disagreement exposes bugs or step-size problems.
\item Differentiating a discretised solver gives the derivative of the solver's
output, which differs from the derivative of the exact solution.
\item Detaching a tensor silently breaks the computational graph; tests must check
gradients, not only values.
\item Relaxing bases to probability vectors makes sequence scores differentiable,
but the optimum of the relaxed problem may not correspond to any discrete
sequence.
\item A design loop proposes candidates; only independent physical evidence
validates them.
\end{itemize}
\end{KeyIdeas}

\section{Twenty worked exercises}
\begin{enumerate}
\item Give $\partial_{k_f}A(0)$. \textbf{Answer:} zero for rate-independent initial state.
\item Differentiate $Ck_r/\lambda$ with respect to $k_f$.
\textbf{Answer:} $-Ck_r/\lambda^2$.
\item Why is there a second term in $s_f$? \textbf{Answer:} the relaxation speed $\lambda$ changes too.
\item State the fixture's local descent direction. \textbf{Answer:} decrease $k_f$ since $dJ/dk_f>0$.
\item Units of $dJ/dk_f$? \textbf{Answer:} population-squared seconds for this loss and rate convention.
\item Why independent scalar FD? \textbf{Answer:} reusing a wrong backward routine cannot validate it.
\item Does smaller $\epsilon$ always help? \textbf{Answer:} no; the $10^{-8}$ error here exceeds $10^{-6}$.
\item Derive the Euler sensitivity update. \textbf{Answer:} differentiate $(1-\lambda h)A_j+h k_rC$ to get $(1-\lambda h)s_j-hA_j$.
\item Which $A$ enters that update? \textbf{Answer:} the old $A_j$, before the simultaneous update.
\item Does correct AD of Euler equal the exact ODE derivative? \textbf{Answer:} not at finite $h$.
\item What catches repeated detach? \textbf{Answer:} full gradient comparison; forward equality alone passes.
\item Add rate units to softplus. \textbf{Answer:} $k=k_0\operatorname{softplus}(\rho)$ with $k_0$ inverse time.
\item Give an unobservable single-time direction. \textbf{Answer:} $(s_r,-s_f)$.
\item Why two times? \textbf{Answer:} their sensitivity rows can span two rate directions; inspect rank and conditioning.
\item Can equilibrium alone infer both absolute rates? \textbf{Answer:} no; proportional scaling leaves equilibrium fixed.
\item Uniform motif score? \textbf{Answer:} $.25$.
\item Uniform target logit gradient for $L=4,\tau=1$? \textbf{Answer:} $.046875$; alternatives $-.015625$.
\item Why do row gradients sum to zero? \textbf{Answer:} softmax is invariant to common row shifts.
\item Show expected GC does not enforce candidate GC. \textbf{Answer:} uniform length-four rows give expected two but A-first argmax gives zero.
\item Design a release gate. \textbf{Rubric:} exact model/units, independent gradient oracle,
solver refinement, parameter domain, identifiability, discrete constraints,
unused validation, documented selection and physical evidence where required.
\end{enumerate}

\section{Bridge to genomic organization}
Chapter 30 distinguishes inherited sequence, regulation, cellular state and
phenotype. A differentiable representation is useful only after those variables,
timescales and observation mappings are explicit. Optimizing a numerical sequence
is not biological evolution, and a smooth model is not a proof that a biological
mechanism has been captured.
